% BHU PYQ -- Transcribed & Corrected
% Paper: ESMV31: Medical Geology
% Exam: B.Sc. (Hons.) Semester III Examination, 2025-26 (NEP)
% Department: Department of Geology
% Paper Ref: CECONF/N/2025-26/050

\documentclass[11pt,a4paper]{article}
\usepackage[a4paper,top=2.5cm,bottom=2.5cm,left=2.5cm,right=2.5cm,headheight=14pt]{geometry}
\usepackage{amsmath,amssymb}
\usepackage[T1]{fontenc}
\usepackage[utf8]{inputenc}
\usepackage{lmodern,microtype,fancyhdr,enumitem,hyperref,lastpage,parskip}

\hypersetup{colorlinks=true,urlcolor=black,linkcolor=black,pdftitle={ESMV31: Medical Geology -- BHU Sem III 2025-26 (NEP VAC)}}
\pagestyle{fancy}\fancyhf{}
\fancyhead[L]{\small\textit{Banaras Hindu University}}
\fancyhead[R]{\small\textit{ESMV31: Medical Geology (2025-26)}}
\fancyfoot[C]{\small Page \thepage\ of \pageref{LastPage}}

\newlist{parts}{enumerate}{2}
\setlist[parts,1]{label=(\alph*),leftmargin=*,itemsep=4pt,topsep=2pt}

\begin{document}

\begin{center}
  {\large\bfseries BANARAS HINDU UNIVERSITY}\\[2pt]
  {\normalsize B.Sc. (Hons.) Semester III Examination 2025-26 (NEP)}\\[6pt]
  \rule{0.8\linewidth}{0.5pt}\\[6pt]
  {\large\bfseries DEPARTMENT OF GEOLOGY}\\[4pt]
  {\large\bfseries Value Added Course (VAC)}\\[4pt]
  {\large\bfseries Paper Code: ESMV31 --- Medical Geology}\\[4pt]
  {\normalsize Time: 2:30 Hours\hfill Maximum Marks: 70}\\[2pt]
  {\small Ref Code: CECONF/N/2025-26/050}
\end{center}

\rule{\linewidth}{0.4pt}\smallskip
\noindent\textit{Instruction: Write your Roll No. at the top immediately on the receipt of this question paper.}\\[2pt]
\noindent\textit{Note: Answer five questions in all, selecting at least two questions each from Section--A and Section--B. Section--C is compulsory. Figures in the right-hand margin indicate marks.}
\bigskip

\section*{SECTION -- A}

\noindent\textbf{Question 1} \hfill [14 Marks]\\
What is a mineral nutrient? Explain in detail why mineral elements are needed for good human health.

\bigskip

\noindent\textbf{Question 2} \hfill [14 Marks]\\
Define the hydrological cycle and discuss in detail the various components of the hydrological cycle with a neat sketch.

\bigskip

\noindent\textbf{Question 3} \hfill [14 Marks]\\
Discuss the biological response of elements with appropriate examples.

\bigskip
\section*{SECTION -- B}

\noindent\textbf{Question 4} \hfill [14 Marks]\\
Write short notes on any \textbf{two} of the following ($2 \times 7 = 14$ Marks):
\begin{parts}
  \item Environmental epidemiology with reference to geogenic factors and disease mapping
  \item Toxicology risk management
  \item Geophagy in humans, its major causes, potential benefits, and associated health risks
\end{parts}

\bigskip

\noindent\textbf{Question 5} \hfill [14 Marks]\\
Write short notes on any \textbf{two} of the following ($2 \times 7 = 14$ Marks):
\begin{parts}
  \item Uptake and regulation of iron in the human body
  \item Drawbacks and health impacts of arsenic consumption on human health
  \item Application of geochemistry to environmental health issues
\end{parts}

\bigskip

\noindent\textbf{Question 6} \hfill [14 Marks]\\
Write short notes on any \textbf{two} of the following ($2 \times 7 = 14$ Marks):
\begin{parts}
  \item Biological functions of major elements
  \item How fluoride overdose in drinking water harms human health
  \item Diversity in natural distribution and abundance of elements
\end{parts}

\bigskip
\section*{SECTION -- C (Compulsory)}

\noindent\textbf{Question 7} \hfill [14 Marks]\\
Answer all the following questions:
\begin{parts}
  \item Write the names of two minerals present in human bone. \hfill [2 Marks]
  \item What are the four major elements of the human body? \hfill [2 Marks]
  \item What is nitrogen fixation? \hfill [1 Mark]
  \item Define Total Dissolved Solids (TDS) and Electrical Conductivity (EC) parameters of water. \hfill [2 Marks]
  \item What is the permissible limit of nitrate for drinking water? \hfill [1 Mark]
  \item What are the carcinogenic associations of fibrous minerals (e.g., asbestos) in the human body? \hfill [3 Marks]
  \item Write the geochemical associations of elements in major igneous and sedimentary rock types. \hfill [3 Marks]
\end{parts}

\end{document}
